% dossier.tex — source dossier on the word Energiens Bestaaen.
% GENERATED by dossier.py; do not edit by hand. Edit selection.txt or
% intro.texfrag beside it instead, then run `make` (see texts/dossiers/README.md).
% This file was produced by the equivalent of:
%   python3 dossier.py 'NEAR(Energi* Bestaa*, 4) OR NEAR(Kraft* Bestaa*, 4) OR NEAR(Stof* Bestaa*, 4) OR NEAR(Materie* Bestaa*, 4) OR NEAR(Energi* Bevar*, 4) OR NEAR(Kraft* Bevar*, 4) OR NEAR(Energi* Konstan*, 4) OR NEAR(Kraft* Konstan*, 4) OR NEAR(Bevægelse* Konstan*, 4) OR NEAR(fysisk* Energi*, 2) OR NEAR(Kraft* Uforgængelig*, 4) OR NEAR(Energi* Omsætning*, 4) OR Ækvivalen* OR Energiform* OR Energisætning* OR Energiprincip* OR Kraftomsætning* OR Energiomsætning* OR Mayer* OR Colding* OR Helmholtz* OR Joule* OR Erhaltung*' --authors hoeffding --title 'Energiens Bestaaen' --highlight 'energi kraftens kraftmængde bestaaen konstans bevarelse bevares ækvivalen uforgængelig kraftomsætning mayer colding helmholtz joule erhaltung' -o <this file> --select texts/dossiers/energibevarelse/selection.txt
% Each passage is followed by a visible \dcite{...} and by sks-search's own
% citation string as a % comment, in the shape sks.el's `w' key produces.
% Compile with lualatex (Makefile) or xelatex.
\documentclass[11pt]{article}
\usepackage[letterpaper,margin=1in]{geometry}
\usepackage{fontspec}
\IfFileExists{libertinus.sty}{\usepackage{libertinus}}{\setmainfont{DejaVu Serif}}
\IfFileExists{danish.ldf}{\usepackage[english,danish]{babel}}{\usepackage[english]{babel}}
\usepackage{microtype}
\usepackage{xcolor,booktabs,longtable}
\usepackage[hidelinks,bookmarksnumbered]{hyperref}
\setcounter{secnumdepth}{0}
\setcounter{tocdepth}{2}
\makeatletter\renewcommand{\@pnumwidth}{2.2em}\makeatother
\emergencystretch=5em
\tolerance=1500
\newcommand{\dcite}[1]{{\nopagebreak\raggedleft\small #1\par}\medskip}
\newcommand{\tlg}[1]{{\setlength{\fboxsep}{0.8pt}\colorbox{yellow}{#1}}}
\IfFontExistsTF{SBL Hebrew}{\newfontfamily\hebrewfont{SBL Hebrew}}{%
  \IfFontExistsTF{Arial Hebrew}{\newfontfamily\hebrewfont{Arial Hebrew}}{%
    \newfontfamily\hebrewfont{DejaVu Sans}}}
\ifdefined\XeTeXversion \TeXXeTstate=1
  \newcommand{\hebr}[1]{{\hebrewfont\beginR #1\endR}}
\else
  \newcommand{\hebr}[1]{{\hebrewfont\textdir TRT #1}}
\fi
\title{\emph{Energiens Bestaaen} hos Høffding\\[0.5ex]\large Udvalgt kildedossier}
\author{Danish Philosophical Texts}
\date{27 September 2026}
\begin{document}
\maketitle

\begin{otherlanguage}{english}
\section{About this dossier}

\noindent A selection from the passages that match this query:

{\raggedright\small\ttfamily NEAR(Energi* Bestaa*, 4) OR NEAR(Kraft* Bestaa*, 4) OR NEAR(Stof* Bestaa*, 4) OR NEAR(Materie* Bestaa*, 4) OR NEAR(Energi* Bevar*, 4) OR NEAR(Kraft* Bevar*, 4) OR NEAR(Energi* Konstan*, 4) OR NEAR(Kraft* Konstan*, 4) OR NEAR(Bevægelse* Konstan*, 4) OR NEAR(fysisk* Energi*, 2) OR NEAR(Kraft* Uforgængelig*, 4) OR NEAR(Energi* Omsætning*, 4) OR Ækvivalen* OR Energiform* OR Energisætning* OR Energiprincip* OR Kraftomsætning* OR Energiomsætning* OR Mayer* OR Colding* OR Helmholtz* OR Joule* OR Erhaltung*\par}

\noindent Bare terms are prefix-matched, so compounds are included; \texttt{NEAR(a b, n)} asks for \texttt{a} and \texttt{b} within \texttt{n} words of each other. They are drawn from in one body of text:

\begin{enumerate}
\item \textbf{Høffding:} 34 passages from 4 works transcribed in this collection. Diplomatic transcriptions; page references are the printed pages recorded in the transcription, and ``p.~109+'' means the passage begins on p.~109 and runs on.
\end{enumerate}

\subsection{Selection}

\noindent Selected for relevance to the law of the conservation of energy, which Høffding calls Energiens (earlier Kraftens) Bestaaen: its forerunners in Descartes' constancy of motion, Leibniz's conservation of force and Lavoisier's conservation of matter; its discovery by Mayer, Colding, Joule and Helmholtz, and their taking it for a truth of reason; its epistemological standing (postulate or experimental law, provable only for closed systems; causal equivalence, Hume's problem, Carnot's principle); its bearing on mind and body (the case against interaction, the identity hypothesis, materialism); and Høffding's analogy between it and the conservation of value (Værdiens Bestaaen). Included with it are passages where psychic energy is measured against the physical law. Omitted: psychic Energi and logical Ækvivalens with no reference to the physical law; Bestaaen in other senses; Helmholtz on sensation and optics; passages in e-books that duplicate a transcription. The selection was made by reading, and is a judgment to be checked, not a mechanical filter.

\noindent Caveats:

\begin{enumerate}
\item The passages were found lexically and then selected by hand (34 kept); see the selection file named in this file's header. Every word form containing \emph{energi}, \emph{kraftens}, \emph{kraftmængde}, \emph{bestaaen}, \emph{konstans}, \emph{bevarelse}, \emph{bevares}, \emph{ækvivalen}, \emph{uforgængelig}, \emph{kraftomsætning}, \emph{mayer}, \emph{colding}, \emph{helmholtz}, \emph{joule} or \emph{erhaltung} is \tlg{highlighted} (167 in all).
\item The quotations are machine-assembled. An entry without a page break of its own carries the page it begins on; for a quotation taken from late in a long entry, check the printed page.
\item 108 further selected passages in e-book reissues of Høffding are not reproduced: those editions are in copyright.
\end{enumerate}

\subsection{Contents by work}
\begin{longtable}{@{}p{0.17\textwidth}p{0.66\textwidth}r@{}}
\toprule
Author & Text & Passages\\
\midrule
\endhead
Høffding & Religionsfilosofi (1901) & 2 \\
Høffding & Den menneskelige Tanke (1910) & 28 \\
Høffding & Religion og Videnskab (1910) & 1 \\
Høffding & Relation som Kategori (1921) & 3 \\
\midrule
& Total & 34\\
\bottomrule
\end{longtable}

\end{otherlanguage}

\tableofcontents
\clearpage

\part*{I.\quad Harald Høffding}
\addcontentsline{toc}{part}{I.\quad Harald Høffding}

\subsection{\textit{Religionsfilosofi} (1901)}

\begin{quote}
I den protestantiske Verden er det temmeligt almindeligt anerkendt, at
der maa skelnes mellem Viden og Tro. Den katolske Kirke anerkender
ikke denne Distinktion i dens fulde Udstrækning og anerkender derfor
heller ikke noget religiøst Problem paa det intellektuelle Omraade.
Den hævder, at Guds Tilværelse kan bevises ad de Veje, Thomas
Aqvinas, støttet til Aristoteles, paaviste, og den har endog erklæret
det for en Trossag at antage, at disse Beviser slaa til. Traditionen
fra Middelalderens store Tid, da “man ikke tvivlede og derfor var
nærmere ved Sandheden” (for at bruge en Sætning af én af Thomas
Aqvinas' moderne Biografer), holdes stadigt i Hævd, og det
anerkendes ikke, at der virkeligt er kommen en Disharmoni ind i det
aandelige Liv. Distinktionen mellem Viden og Tro indrømmer, at en
saadan Disharmoni er til Stede, — at Aandslivets Enhed er brudt. Er
nu denne Distinktion selv kun en Station paa Vejen, idet Udviklingen
vil føre til en Udskillelse af al Religion af Aandslivet, fordi kun
det kan bestaa, der kan træde i harmonisk Forhold til Livets andre
Elementer? — Hvad enten man nu antager, at det kun er Religionens
klassiske Tid, der er forbi, eller at dens Tid overhovedet er forbi,
saa kan det Spørgsmaal ikke skydes til Side, om der gaar aandelig
Værdi tabt ved den Opløsningsproces, der er indtraadt, eller om der
kan paavises Erstatning for den aandelige Enhed og Sluttethed, som
tabtes. Allerede denne Betragtning viser, at Spørgsmaalet om
Værdiens \tlg{Bestaaen} er en væsentlig Side ved det religiøse Problem,
hvilken Løsning af dette man saa end helder til. Det er ikke mindst
fra denne Side, at Filosofien i det sidste halvandet Aarhundrede (fra
Rousseau's, Lessing's og Kant's Tid af) har behandlet det religiøse
Problem. Baade Romantikere og Positivister have været opmærksomme paa
denne Side ved det. — Det maa jo vel erindres, at selvom man kunde
paavise, at der ingen aandelig \tlg{Energi} gik tabt ved Religionens
Opløsning, var det ikke sagt, at \tlg{Energien} under de nye Forhold fik en
ligesaa værdifuld Anvendelse som i Religionens klassiske Tider.
Ligesaa lidt paa det aandelige som paa det materielle Omraade er
\tlg{Energiens} \tlg{Bestaaen} uden videre det Samme som Værdiens \tlg{Bestaaen}.
\tlg{Energien} kan bestaa, uden at Værdien bestaar, medens det Omvendte
ikke vel kan tænkes, da \tlg{Energi} i og for sig vil være en Betingelse
for det Værdifuldes \tlg{Bestaaen} og Udvikling. Mange Fritænkere gaa uden
videre ud fra, at det menneskelige Liv vil blive kraftigere og
rigere, naar Religionen falder bort, — en Antagelse, som ikke er
selvfølgelig, og som hviler paa Forudsætningen om, at der er psykiske
\tlg{Ækvivalenter} til Stede, — \tlg{Ækvivalenter} baade hvad Værdi og hvad
\tlg{Energi} angaar. Det vil da blive Opgaven at paavise disse \tlg{Ækvivalenter}
og saaledes godtgøre Værdiens \tlg{Bestaaen}. Men det er et stort
Spørgsmaal og en væsentlig Side ved det religiøse Problem, om
saadanne \tlg{Ækvivalenter} kunne paavises. — Iøvrigt vil Problemet om
Værdiens \tlg{Bestaaen} ogsaa opstaa indenfor Kirken, ti Religionen og det
religiøse Liv undergaar stadige Forandringer; og Spørgsmaalet bliver,
om Værdierne fra Religionens klassiske Tider holde sig gennem disse
Forandringer, f. Ex., om Urkristendommen, baade hvad Værdi og hvad
\tlg{Energi} angaar, er bevaret i vore Dages Kristendom. I sin skarpeste
Form er dette Problem fremtraadt i den Paastand: “det nye
Testamentes Kristendom er ikke til.” —
\end{quote}
\dcite{Høffding, Religionsfilosofi, pp. 8–10  [hoeffding/religionsfilosofi:16]}
% hoeffding/religionsfilosofi:16  (Religionsfilosofi; pp. 8–10)

\begin{quote}
Dersom Sætningen om Værdiens \tlg{Bestaaen} skulde drøftes i hele dens
Rækkevidde, vilde Diskussionen blive uafsluttelig, og der vilde
desuden snart mangle klare Udgangspunkter for at føre den. Jeg
fremhæver derfor straks de Grænser, indenfor hvilke min Undersøgelse
vil bevæge sig, og som maa haves i Erindring i alt Følgende. — Naar
jeg forsøger paa at vise, at den nævnte Sætning er Grundtanken i al
Religion, vil dette Forsøg ikke være mislykket, fordi det viser sig,
at ingen Religion udtrykker Sætningen med Klarhed og Konsekvens. Det
vil være tilstrækkeligt, om der kan paavises en udtrykkelig Trang
eller Tendens til at fastholde denne Grundtanke, saaledes at
Maalestokken ved Vurderingen af Religionerne i deres indbyrdes
Forhold bliver den Grad, i hvilken de formaa at udtale og gennemføre
Sætningen. — Naar der tales om Værdiens \tlg{Bestaaen}, saa kunde dette
forstaas saaledes, at Værdien ikke forsvinder, men at der stadigt er
Værdi i Tilværelsen, hvad enten den aftager eller voxer eller skifter
mellem Aftagen og Tiltagen. Men jeg opfatter Udtrykket Værdiens
\tlg{Bestaaen} i Analogi med Udtrykket \tlg{Energiens} \tlg{Bestaaen}, saaledes at
Sætningen hævder en stadig \tlg{Bestaaen} af Værdien gennem alle
Formforandringer. Det vil maaske herved blive nødvendigt at skelne
mellem potentiel og aktuel Værdi, ligesom man skelner mellem potentiel
og aktuel \tlg{Energi}, og den ene Distinktion er ligesaa tydelig og
ligesaa berettiget som den anden. — Men Værdiens \tlg{Bestaaen} gennem
alle Formforandringer og trods Forskellen mellem Mulighed og
Virkelighed er dog kun et Minimum. Der vil kunne paastaas, at den
blotte \tlg{Bestaaen} af Værdien ikke er tilfredsstillende og egentligt
indeholder en Modsigelse, da Værdien kun kan bestaa ved at vokse, da
blot Gentagelse sløver, medens paa den anden Side Forandring kun er
værdifuld, naar den bringer Forhøjelse. At en Forøgelse af Værdien
ikke strider mod \tlg{Energiens} \tlg{Bestaaen}, behøver ingen nærmere
Paavisning; Værdiforøgelsen forudsætter blot en ny og fuldkomnere
Anvendelse af \tlg{Energien}, ikke en Forøgelse af denne i det Hele. Det
er for Simpelhedens Skyld, at jeg begrænser Undersøgelsen til
Værdiens \tlg{Bestaaen}. Og det vil tillige vise sig, at den religiøse
Bevidsthed i Grunden ogsaa har begrænset sig hertil. Værdiens
\tlg{Bestaaen} vil stille os Vanskeligheder og Opgaver nok, selvom vi ikke
indlade os paa Alt, hvad der ligger i den uundgaaelige Konsekvens, at
Værdien kun kan bestaa, dersom den forøges. Dog maa vi ikke glemme,
at denne Konsekvens er til Stede. — Endeligt kunde man, naar der er
Tale om Værdiens \tlg{Bestaaen}, spørge: hvilken Værdi og hvor megen Værdi?
Paa disse Spørgsmaal kommer jeg kun ind, forsaavidt Forskellen mellem
de forskellige Religioner og religiøse Standpunkter væsentligt beror
paa de forskellige Svar, de give paa dem. Det er klart, at Sætningen
om Værdiens \tlg{Bestaaen} kan antages, hvilke Svar man saa giver.
Hypotesen om, at Værdiens \tlg{Bestaaen} er den religiøse Grundsætning,
rokkes ikke ved, at de forskellige Religioner og religiøse
Standpunkter afvige i høj Grad fra hverandre i Henseende til
Beskaffenheden og Omfanget af det Værdifulde, paa hvis \tlg{Bestaaen} de
tro.
\end{quote}
\dcite{Høffding, Religionsfilosofi, pp. 10–12  [hoeffding/religionsfilosofi:19]}
% hoeffding/religionsfilosofi:19  (Religionsfilosofi; pp. 10–12)

\subsection{\textit{Den menneskelige Tanke} (1910)}

\begin{quote}
Den psykiske \tlg{Energi}, der ytrer sig i Syntesen, maale vi ved at tage
Hensyn til alle disse Omstændigheder. Om exakt Maaling er der her
naturligvis ikke Tale, kun om et vurderende Skøn. Men saasnart vi søge
at begrunde, hvorfor et Tankearbejde, en Aandsretning, en Livsanskuelse
kan kaldes lavere eller højere end en anden, maa vi tage Hensyn til
alle disse Omstændigheder, og vi gøre det i Virkeligheden ogsaa, mere
eller mindre bevidst, og mere eller mindre nøjagtigt. Der er mindre
psykisk \tlg{Energi} til Raadighed i Drømmetilstanden, hvor de skiftende
Emner udløse hvert sit Forestillingsløb, end i den vaagne Tilstand,
hvor herskende Formaalstanker medføre en Hovedretning. Det er en rent
psykologisk Maalestok, vi her anvende. Og selvom vi ved Dannelsen af
Begrebet psykisk \tlg{Energi} have havt Vejledning i Analogien med den
fysiske \tlg{Energis} Begreb, saa er hint Begreb dog bygget paa sine
særegne Erfaringer, ved hvilke fysiske eller fysiologiske Iagttagelser
og Undersøgelser ikke behøve at spille nogen afgørende Rolle. Derfor
behøve vi ikke at miskende Betydningen af de Forsøg, der fra
experimental-psykologisk og fysiologisk Side ere gjorte paa at maale
det Hjernearbejde, der forudsættes ved visse elementære
Aandsvirksomheder, f. Ex. Regnen eller Husken, gennem dets Forhold
til det Muskelarbejde, der kan øves samtidigt dermed. Man kan mærke
saadant Tankearbejde ved den Forminskelse, det fremkalder i det
samtidigt udførte Muskelarbejde. Men selve denne Maaling forudsætter,
at man ad rent psykologisk Vej kan skelne mellem vanskeligere og
lettere Aandsarbejde; ellers bevæger man sig i en Kreds. Det
Hjernearbejde, der svarer til det i det Foregaaende beskrevne
aandelige Arbejde, kender man jo ikke direkte, kun gennem selve det
Aandsarbejde, til hvilket det antages at svare.
\end{quote}
\dcite{Høffding, Den menneskelige Tanke, pp. 6–7  [hoeffding/menneskelige-tanke:19]}
% hoeffding/menneskelige-tanke:19  (Den menneskelige Tanke; pp. 6–7)

\begin{quote}
7. Begrebet \tlg{Energi} i den Betydning, i hvilken vi her have taget det,
har sit Udspring i Naturvidenskaben, og det er ad Analogiens Vej, vi
have forsøgt at overføre det paa Sjælelivets Omraade. Alle
psykologiske Udtryk ere jo overførte fra det Materielle til det
Aandelige, fra Rummets Verden til Aandens og Tankens Verden. Men der
er Vanskeligheder ved Gennemførslen af Analogien. En exakt Maaling er
paa det psykiske Omraade kun mulig, hvad forholdsvis elementære
sjælelige Fænomener angaar og aldrig paa saa fuldkommen Maade som for
den materielle \tlg{Energis} Vedkommende. Den moderne Fysiologi staar, hvad
Metode angaar, paa det Standpunkt, at den kræver fysiologiske
Tilstande forklarede af fysiologiske Aarsager, forsaavidt ikke fysiske
Aarsager fra Omverdenen gribe ind. Enhver Hjernetilstand eller
Hjernevirksomhed skal derfor ogsaa forklares ud fra Hjernens og
Organismens tidligere Tilstande. Efter strengt naturvidenskabelig
Metode maa enhver materiel Tilstand i og udenfor Organismen finde sin
Forklaring ved at ses som opstaaet ved Omsætning i ny Form af den
fysiske \tlg{Energi}, der ytrede sig i forudgaaende materielle Tilstande.
Naar to Hjernetilstande altsaa afløse hinanden, vil derfor Forklaringen
først være funden, naar det paavises, at al den \tlg{Energi}, der ytrede sig
i den første Tilstand, uden Rest er bleven omsat i den anden Tilstand,
forsaavidt den ikke er straalet ud i andre fysiologiske Processer. Det
er maaske umuligt at gennemføre Experimenter i denne Retning; men
saadanne vilde betegne Fuldendelsen af den fysiologiske Metode, der
som sit Motto kan sætte Spinozas Ord, at en Forklaring af en
legemlig Tilstand ved Sjælens Indgriben er Et med en Indrømmelse af,
at man ikke kan forklare den. Selvom nu et af de Forsøg, man kunde
gøre, syntes at vise, at der ved en organisk Tilstands Opstaaen opstod
eller forsvandt fysisk \tlg{Energi} uden fysisk \tlg{Ækvivalent}, vilde man fra et
rent naturvidenskabeligt Standpunkt snarere tro, at der var en Fejl
ved Experimentet, end slaa sig til Ro ved, at man havde opdaget en
Grænse for den fysiske \tlg{Energis} \tlg{Bestaaen}. „Principet om \tlg{Energiens}
\tlg{Bestaaen},“ siger Maxwell [note: Scientific Papers. II.
p. 759.], „har faaet saa stor videnskabelig Vægt, at ingen Fysiolog
vilde have nogen Tillid til et Forsøg, der viste en betydelig Forskel
mellem det af et levende Væsen udførte Arbejde og Summen af modtagen
og forbrugt \tlg{Energi}.“ Siden Maxwell nedskrev denne Sætning, er der
foretaget Forsøg, der — indenfor visse Fejlgrænser — vise, at
Stof- og Kraftskiftet under Udførelse af aandeligt og legemligt
Arbejde frembyder \tlg{Ækvivalens} mellem, hvad Organismen optog, og hvad
den forbrugte [note: Atwater: Neue Versuche über Stoff- und
\end{quote}
\dcite{Høffding, Den menneskelige Tanke, pp. 22–23  [hoeffding/menneskelige-tanke:51]}
% hoeffding/menneskelige-tanke:51  (Den menneskelige Tanke; pp. 22–23)

\begin{quote}
Kraftwechsel des menschlichen Kørpers (Ergebnisse der Physiologie
III, 1. 1904).]. Det vil herefter blive endnu vanskeligere end før at
antage, at der i nogen Del af Organismen skulde foregaa Noget, der
strider mod \tlg{Energisætningen}.
\end{quote}
\dcite{Høffding, Den menneskelige Tanke, p. 23  [hoeffding/menneskelige-tanke:52]}
% hoeffding/menneskelige-tanke:52  (Den menneskelige Tanke; p. 23)

\begin{quote}
Men hvorledes vil det da under disse Forhold blive muligt at fastholde
den psykiske \tlg{Energi} som en særegen Form for \tlg{Energi}? Er det da ikke
rimeligere at antage, at al \tlg{Energi} er fysisk, og at den psykiske
\tlg{Energi}, vi i det Foregaaende have talt om, kun er et Symptom paa en
fysisk \tlg{Energi}, vi ingen direkte Adgang have til? — Vi kunne jo ikke
gennemføre \tlg{Ækvivalensprincipet} paa selve det psykiske Omraade,
saaledes som paa det fysiske Omraade. Der sker Afbrydelser indenfor et
og samme individuelt Bevidsthedsliv (ved Søvn, Afmagt o. s. v.), og
imellem forskellige individuelle Bevidstheder kunne vi ikke paavise
psykisk Kontinuitet, saaledes som vi mellem forskellige Legemer,
f. Ex. mellem to Organismer, kunne paavise fysisk Kontinuitet.
\end{quote}
\dcite{Høffding, Den menneskelige Tanke, p. 24  [hoeffding/menneskelige-tanke:53]}
% hoeffding/menneskelige-tanke:53  (Den menneskelige Tanke; p. 24)

\begin{quote}
Hvor der trods Alt viser sig en Mangel paa Sammenhæng paa det psykiske
Omraade, behøve vi dog ikke at opgive Kontinuitetens Princip, ligesaa
lidt som Fysikeren opgiver det, naar \tlg{Energien} synes at være forsvunden
i en Hvile- eller Ligevægtstilstand. Ligesom Fysikeren indfører
Begrebet potentiel fysisk \tlg{Energi} for at udtrykke den Kendsgerning, at
hvad der syntes at være forsvundet, kan udløses igen, saaledes vil
Psykologen være berettiget til at tale om potentiel psykisk \tlg{Energi},
især støttende sig til den Kendsgerning, at psykiske Elementer kunne
reproduceres efter et Mellemrum, i hvilket de ikke have været i
Bevidstheden. Ord som „Anlæg“, „Mulighed“, „Disposition“, „Evne“,
„Spor“, „Tendens“ udtrykke det Samme som Ordet potentiel \tlg{Energi}.
Psykologen er dog ved Anvendelsen af dette Begreb uheldigere stillet
end Fysikeren ved Anvendelsen af sit tilsvarende Begreb. For det
første kan Fysikeren ligefrem pege paa det Fænomen, i hvilket den
potentielle \tlg{Energi} har sit Sæde, f. Ex. den spændte Fjeder eller den
hvilende Sten paa Taget. Derimod betyder paa det psykiske Omraade
fuldstændig Ligevægt (Hvile) det Samme som Ubevidsthed, altsaa noget
for Iagttagelsen Utilgængeligt (hvad det saa ellers i sig selv maatte
være). Og for det Andet kan man paa det fysiske Omraade foretage
Omsætninger, der vise, at der kan udløses et Kvantum af \tlg{Energi}, som
svarer til det Kvantum, der tilsyneladende forsvandt ved Overgang til
den „potentielle“ Tilstand. Men paa det psykiske Omraade kan det ikke
paa denne Maade godtgøres, at Anlæg og Udvikling svare til hinanden.
Naar Begrebet potentiel psykisk \tlg{Energi} alligevel har sin Berettigelse,
er det, fordi det staar som Udtryk for en uløst, men uafviselig
Opgave, og fordi den indeholder en Opfordring til stedse at gaa videre
i Iagttagelse og Analyse for at finde saa megen Kontinuitet som muligt
mellem Sjælelivets forskellige Tilstande. Der er en lang Vej til
absolut Hvile — og det kunde jo være, at den ligesaalidt forekommer
paa det psykiske som paa det fysiske Omraade. Mulighed bestaar overalt
i en Virkelighed, hvad enten denne er os tilgængelig eller ikke.
\end{quote}
\dcite{Høffding, Den menneskelige Tanke, pp. 30–31  [hoeffding/menneskelige-tanke:62]}
% hoeffding/menneskelige-tanke:62  (Den menneskelige Tanke; pp. 30–31)

\begin{quote}
Det er den naturvidenskabelige Metode, der Skridt for Skridt fører til
en saadan Opfattelse. Inertisætningen kræver, at enhver Forandring i
et materielt Punkts Tilstand (af Hvile eller retliniet Bevægelse med
en vis Hastighed) skal udledes af et andet materielt Punkts
Indflydelse. Dermed er det udelukket, at en ikke-materiel (d. v. s.
ikke-rumlig) Aarsag kan faa naturvidenskabelig Betydning. Derved er det
Princip fastslaaet, at et materielt Fænomen skal have en materiel
Aarsag. Maxwell finder Inertisætningens Erfaringsbevis deri, at
vi, hvergang vi træffe en Forandring i et Legemes Bevægelsestilstand,
kunne føre denne Forandring tilbage til en Virken mellem dette Legeme
og et andet Legeme. Dette Erfaringsbevis er Naturvidenskaben stadigt
ifærd med at føre. — Til Inertisætningen slutter sig Sætningen om
den fysiske \tlg{Energis} \tlg{Bestaaen}, der som allerede omtalt synes at maatte
fastholdes ogsaa for den organiske Naturs Vedkommende.
\end{quote}
\dcite{Høffding, Den menneskelige Tanke, p. 31+  [hoeffding/menneskelige-tanke:64]}
% hoeffding/menneskelige-tanke:64  (Den menneskelige Tanke; p. 31+)

\begin{quote}
Den Opfattelse, vi her slutte os til, vender egentligt tilbage til den
barnlige og praktiske Anskuelse, der betragter Legemet som det ydre,
for Sanserne tilgængelige Udtryk for Personligheden. Og vi faa derved
tillige den bedste og mest holdbare Arbejdshypotese, idet vi opfordres
til paa en Gang at følge de psykologiske og de fysiologiske
Fænomeners Rækker hver for sig saa langt som overhovedet muligt, og at
undersøge Forholdet mellem dem, Analogierne og Proportionaliteterne,
saa nøje som muligt. Det er et Funktionsforhold i matematisk Betydning
af Ordet, denne Opfattelse antager mellem de sjælelige og de materielle Fænomener: vi søge, og ofte finde vi ogsaa, at en Varieren af de
sjælelige Tilstande svarer til en Varieren af de legemlige Tilstande,
særligt Hjernetilstandene, og omvendt. Og hertil indskrænker sig i
Virkeligheden den Viden, vi have om disse Ting. Vor Hypotese (for
hvilken det mest passende Navn er Identitetshypotesen, da den ikke
antager, at Sjæl og Legeme ere to forskellige Væsener) betragte vi
foreløbigt kun som en Arbejdshypotese. Om den giver en afsluttende
Forklaring af Forholdet mellem det Aandelige og det Materielle, mellem
psykisk og fysisk \tlg{Energi}, er et andet Spørgsmaal, som i en senere
Sammenhæng vil blive drøftet. Selv Forskere, der forkaste vor Hypotese
som en definitiv Løsning, indrømme dog, at den er en nyttig
Arbejdshypotese. Saaledes H. Bergson og C. Stumpf. Om
denne Indrømmelse fra deres Standpunkt er konsekvent, er et Spørgsmaal
for sig. En Hypotese, der slet ikke har nogen Grund i Tingenes Natur,
vil ikke i Længden egne sig til Arbejdshypotese. —
\end{quote}
\dcite{Høffding, Den menneskelige Tanke, pp. 32–33  [hoeffding/menneskelige-tanke:68]}
% hoeffding/menneskelige-tanke:68  (Den menneskelige Tanke; pp. 32–33)

\begin{quote}
Det er nu ikke altid, at begge de Emner, der skulle sammenlignes, ere
givne eller uden videre kunne fremkaldes i Erindringen. Ofte er kun det
ene Emne til Raadighed, og Opgaven er at finde et andet Emne, der staar i
et bestemt Ligheds- eller Forskelsforhold til det. Jo klarere og
tydeligere dette Forhold kan forestilles, des større Mulighed er der
for, at det nye Emne kan findes. Den nøjagtigste Angivelse af hint
Forhold have vi, hvor Opgaven kan sættes i Ligning. Vi have da en
bestemt Maalestok til at afgøre, om det søgte Emne optræder eller ikke,
og ved en metodisk Fremgangsmaade kunne vi maaske fremtvinge det Emne,
der fyldestgør Fordringen. Paa mere ubestemt Maade sker noget Lignende,
naar vi ville fremkalde Noget i vor Erindring og kun vide om det, at det
staar i en vis Sammenhæng med en Forestilling, der allerede er fremme i
os; ved Koncentration af Opmærksomheden paa denne, eller som man har
kaldt det, ved en Boring, kan da det Søgte springe frem. Den psykiske
\tlg{Energi} fremtræder ganske særligt i saadanne Tilfælde. Det gælder da ikke
blot om at forbinde givne Led, men at finde et Led, som maaske er helt
nyt, men som staar i et bekendt Forhold til et givet Led. Hvor vi vel
kunne bestemme Forholdet til det givne Led, men paa Grund af vor
Erfarings Begrænsning ikke kunne finde det nye Led, der ende vi i et
Problem. Saaledes ere vi f. Ex. stillede, naar vi antage, at der paa
det psykiske Omraade gives noget Tilsvarende til Forskellen mellem
potentiel og aktuel \tlg{Energi} paa det fysiske Omraade; ti om potentiel
psykisk \tlg{Energi} kunne vi ingen positiv Forestilling danne os; vi kunne
kun sige, at den forholder sig til aktuel psykisk \tlg{Energi} som potentiel
\tlg{Energi} forholder sig til aktuel \tlg{Energi} paa det fysiske Omraade. (Smlgn.
ovenfor 8–9).
\end{quote}
\dcite{Høffding, Den menneskelige Tanke, p. 66  [hoeffding/menneskelige-tanke:150]}
% hoeffding/menneskelige-tanke:150  (Den menneskelige Tanke; p. 66)

\begin{quote}
Formedelst Loven om \tlg{Energiens} \tlg{Bestaaen} kunne de forskellige
Naturkræfter betragtes som \tlg{ækvivalente} og substitueres for hverandre,
hvor det blot gælder selve \tlg{Energisynspunktet}, ikke \tlg{Energiens} specielle
Form eller dens praktiske Betydning. Potentiel \tlg{Energi} er saaledes fra dette Synspunkt identisk med aktuel \tlg{Energi}. Paa det psykiske Omraade
kunne vi indenfor visse Grænser anvende et tilsvarende Synspunkt. Vi
kunne ogsaa her tale om \tlg{Ækvivalenter}. Saaledes er Selvbeherskelse kun
mulig formedelst forskellige Drifters eller Interessers \tlg{Ækvivalens}. En
Drift kan (naar den ikke dør en langsom Død af Mangel paa Næring) ikke
fortrænges uden ved, at en anden Drift drager al \tlg{Energi} til sig, som
naar f. Ex. Kunstinteressen fortrænger Driften til sanselig Nydelse.
Identitetssynspunktet er her: Noget, som opfylder Sindet og lægger
Beslag paa dets \tlg{Energi}.
\end{quote}
\dcite{Høffding, Den menneskelige Tanke, p. 179+  [hoeffding/menneskelige-tanke:437]}
% hoeffding/menneskelige-tanke:437  (Den menneskelige Tanke; p. 179+)

\begin{quote}
91. Er Idealet nu naat, naar Aarsagsforholdet antager Kontinuitetens Form? Maa man
ikke gaa et Skridt videre? — Analogien med Forholdet mellem Grund og Følge (88)
synes at kræve det. Ti i dette Forhold gælder det, at der ikke maa være Noget i
Følgen (Slutsætningen), som ikke var i Grunden (Forudsætningerne), og det
fuldkomneste Begrundelsesforhold kræver endog, at Grund og Følge kunne byttes om, saa at \tlg{Ækvivalensen} gælder i begge Retninger (86). Denne Analogi førte det 17.
Aarhundredes Filosofer (for hvem der her var Mere end Analogi) til at formulere
Aarsagssætningen saaledes, at den krævede, at der ikke maatte være Mere i
Virkningen end i Aarsagen, og at der maatte være absolut Ensartethed mellem Aarsag og
Virkning. Hos Spinoza fremtræder denne Forudsætning skarpest udpræget. Den
kan endnu spores hos Kant, og især hos Nogle af hans Efterfølgere
(Hamilton og Spencer), der gøre Skridtet fuldt ud og finde en
gensidig \tlg{Ækvivalens} indeholdt i Kausalitetsbegrebet. Af stor erkendelsesteoretisk
Interesse er det, at denne Opfattelse spillede en væsentlig Rolle ved Opdagelsen
af Loven om \tlg{Energiens} \tlg{Bestaaen}. Baade Robert \tlg{Mayer}, \tlg{Colding} og
\tlg{Joule} gik udfra, at \tlg{Energiens} \tlg{Bestaaen} var en nødvendig, i selve
Aarsagsbegrebet indeholdt Forudsætning. \tlg{Mayer} blev aldrig ked af at gentage
Sætningen causa æqvat effectum, der for ham var en Følge af „den logiske Grunds
Lov"; det var for ham en selvfølgelig Sætning. \tlg{Colding} fandt den modsatte Antagelse
„fornuftstridig", og \tlg{Joule} kaldte den „absurd". Heller ikke \tlg{Helmholtz} havde
Bevidstheden om at opstille en ny Sætning [note: Smlg. Den nyere Filosofis
Historie2 II. p. 497–499; 597.]. Heldigvis stansede disse
Forskere ikke ved denne Betragtning, men underkastede deres Synsmaade en
experimental Prøvelse.
\end{quote}
\dcite{Høffding, Den menneskelige Tanke, pp. 213–214  [hoeffding/menneskelige-tanke:518]}
% hoeffding/menneskelige-tanke:518  (Den menneskelige Tanke; pp. 213–214)

\begin{quote}
Overfor denne Betragtningsmaade maa det hævdes, at Aarsagsbegrebet ikke vilde
miste al Betydning, fordi det ikke lykkedes at godtgøre \tlg{Ækvivalens} mellem Aarsag
og Virkning. Der kunde gælde et Aarsagsforhold mellem Emnerne, selvom der fandt en
stadig Variation Sted, naar blot Fremtiden var analog med Fortiden paa den Maade,
at vi deraf, at Begivenheden \$B\_1\$ følger efter Begivenheden \$A\_1\$, kunde slutte,
at Begivenheden \$B\_2\$ vil følge efter \$A\_2\$, — naar \$A\_1,A\_2,B\_1,B\_2\$ betegne
\end{quote}
\dcite{Høffding, Den menneskelige Tanke, p. 214  [hoeffding/menneskelige-tanke:519]}
% hoeffding/menneskelige-tanke:519  (Den menneskelige Tanke; p. 214)

\begin{quote}
forskelligartede Begivenheder. Selv hvis de kvalitative Forskelligheder vare irrekdutible, vilde Aarsagssætningen kunne gælde. Saa meget desmere vilde den kunne
gælde, selvom der ikke kunde paavises kvantitativ \tlg{Ækvivalens} mellem Aarsag og
Virkning, saa at f. Ex. A var Aarsag til \$B = A + \$, og B Aarsag til
\$C = B y\$. Lovmæssighed betyder ikke nødvendigvis \tlg{Ækvivalens}; denne staar
ikke som selvfølgelig.
\end{quote}
\dcite{Høffding, Den menneskelige Tanke, p. 214+  [hoeffding/menneskelige-tanke:520]}
% hoeffding/menneskelige-tanke:520  (Den menneskelige Tanke; p. 214+)

\begin{quote}
Det kvantitative Forhold mellem Aarsag og Virkning, som det for mange Emners
Vedkommende er lykkedes at paavise, faar erkendelsesteoretisk set især sin
Betydning ved, at det kun ved dets Hjælp er muligt at anvende de fundne
Aarsagsforhold paa fremtidige Emner. Rent kvalitativt kan nemlig Forholdenes
Identitet ikke med Nøjagtighed bestemmes. Den Grad af Lighed, der maa være tilstede
mellem \$A\_1\$ og \$A\_2\$, for at det skal være berettiget at vente \$B\_2\$ efter \$A\_2\$,
ligessom \$B\_1\$ kom efter \$A\_1\$, kan kun fastsættes ved et Skøn. Kun ved Tælling,
Maaling og Vejning kan Forholdenes Identitet godtgøres, saa at sikker Genkendelse
bliver mulig, — og selv da kun tilnærmelsesvis. Skønt der ikke kan paavises absolut
Gentagelse, hverken i Naturen eller i Historien, er det dog naturligt, at Tanken
søger at naa saa langt i at konstatere Forholdenes Identitet, som det paa nogen
Maade er muligt. Det er da ogsaa herpaa Opmærksomheden først og fremmest er rettet
indenfor de specielle Forskningsgrene, og vi maa tage det som en Kendsgerning, at
jo nøjagtigere Forholdenes Identitet kan godtgøres, des mere viser det sig, at der
er \tlg{Ækvivalens} mellem hvad der gaar forud og hvad der følger efter.
\end{quote}
\dcite{Høffding, Den menneskelige Tanke, pp. 214–215  [hoeffding/menneskelige-tanke:521]}
% hoeffding/menneskelige-tanke:521  (Den menneskelige Tanke; pp. 214–215)

\begin{quote}
Men selvom \tlg{Ækvivalensforholdet} gennemførtes overalt, vilde der dog vedblive at
være et historisk Element i vor Erkendelse, der er givet med den Retning, i hvilken
Omsætningen fra det ene Emne til det andet foregaar. Der kan være kausal \tlg{Ækvivalens}
uden at være \tlg{Værdiækvivalens}. Det er ikke blot den menneskelige Historie, der
viser, at vi i Aarsagernes Række ikke lige saa godt kunne gaa tilbage som frem;
ogsaa Naturen har sin Historie, der fysisk set især karakteriseres ved den
Omstændighed, at Omsætningen fra Varme til Bevægelse ikke foregaar med samme Lethed
som Omsætningen fra Bevægelse til Varme. Den rent kausale Betragtning viser sig her
utilstrækkelig, og der er Brug for andre reale Kategorier end Aarsagskategorien,
der er ligegyldig mod Retningen. Dette gælder, hvad enten Kausalforholdet kan
opfattes som \tlg{Ækvivalens} eller ikke. Og dette gør, at Kontinuitetssynspunktet i
Virkeligheden bliver det højeste Synspunkt, fra hvilket Aarsagsforholdet kan ses.
\end{quote}
\dcite{Høffding, Den menneskelige Tanke, p. 215  [hoeffding/menneskelige-tanke:522]}
% hoeffding/menneskelige-tanke:522  (Den menneskelige Tanke; p. 215)

\begin{quote}
b. Totalitet — 93. Selv hvor vi i Aarsagsrækken kunne gaa baade frem og tilbage, ved gensidig
\tlg{Ækvivalens}, er, som allerede bemærket, Retningen ikke ligegyldig. \tlg{Ækvivalens} i
Henseende til \tlg{Energi} er ikke det eneste Synspunkt. Det viser sig psykologisk deri,
at en sjælelig Tilstand ikke blot er beskreven ved den sammenfattende \tlg{Energi}, den
raader over, men ogsaa karakteriseres ved den kvalitative Beskaffenhed af de
psykiske Emner (Tanker, Interesser, Formaal), der udgøre Tyngdepunktet, det reale
Selv. Derved bestemmes Sjælelivets Retning til enhver Tid, den Stræben, i hvis
Tjeneste Syntesen tilsidst staar. Ti den sammenfattende Virken stiller ikke alle
Elementer paa lige Trin; der er Elementer, som indtage en central Plads, vore
stadige Tanker, vor inderste Trang og Higen, — og der er andre, som indtage en
mere periferisk Plads. Ombytning af disse Pladser kan være Et med en sjælelig
Revolution [note: Smlg. min Afhandling om Begrebet Villie i første Bind af
Tidsskriftet „Psyche".]. Og en Analogi hertil findes paa det materielle Omraade.
Det er ikke nok at undersøge Forholdet mellem Aarsag og Virkning indbyrdes. Dette
gøres bedst, naar man kan isolere en Kreds af Emner fra alle andre. Sætningen om
\tlg{Energiens} \tlg{Bestaaen} kan saaledes kun experimentalt godtgøres, naar man kan danne
isolerede Helheder og holde Indflydelser udefra borte. Men det er da altid tilsidst
en vilkaarlig Begrænsning, en Abstraktion. Og Spørgsmaalet bliver, naar man vender
tilbage til de virkelige Forhold, stedse om, i hvilken Retning Aarsagsrækken under
sit fortsatte Løb bevæger sig, — hvilke Led der have Mulighed for her paany at
gøre sig gældende. Retning hører ligesaa væsentligt med til en Bevægelse som Masse
og Hastighed. Leibniz opstillede derfor, foruden Loven om \tlg{Kraftens} \tlg{Bestaaen}, en Lov om
Retningens \tlg{Bestaaen}. Ifølge Inertisætningen ændrer Bevægelse ikke sin Retning,
saalidt som sin Hastighed, uden ved ydre Kræfters Indflydelse. Men selv naar
Retningen ændres, er den tidligere Retning dog medvirkende ved den nye Retning, der
er Resultanten af den tidligere Retning og den ydre Krafts Retning. Der er altsaa
Kontinuitet i Retningen ligesom i \tlg{Energien}. Retning er ikke Noget, der først dukker
op paa et bestemt Punkt af Bevægelsen, men er med saa langt, vi gaa tilbage, og de
tidligere Retninger bestemme de følgende, ligesom paa det psykiske Omraade vor
tidligere Stræben og Villen bestemmer den senere.
\end{quote}
\dcite{Høffding, Den menneskelige Tanke, pp. 218–219  [hoeffding/menneskelige-tanke:527]}
% hoeffding/menneskelige-tanke:527  (Den menneskelige Tanke; pp. 218–219)

\begin{quote}
Og som enhver Helhed fra et eller andet Synspunkt stedse vil kunne betragtes som
Del, saaledes kan enhver Del igen betragtes som Helhed. Rækken af Dele og Helheder
(Række (5)) er omvendt symmetrisk (en absolut Korrelatrække, se 68). Allerede heri
ligger det, at absolute Atomer ikke kunne tænkes. Der gives ligesaalidt nogen Del,
der ikke kan være Helhed, som der gives nogen Helhed, der ikke kan være Del. Det har
da ogsaa vist sig, at en saadan Antagelse fører til Modsigelse; allerede
Huyghens og Leibniz saa, at den stred mod \tlg{Energiens}
\tlg{Bestaaen} [note: Den nyere Filosofis Historie.2 I. p.
346.], og i den nyeste Tid er det blevet mere og mere klart, at Atomet selv staar
som en hel lille Verden, indenfor hvilken nye Problemer dukke op. Heller ikke paa
det psykiske Omraade kunne vi stanse ved sidste, aldeles usammensatte Elementer;
selv de simpleste Fornemmelser vise sig, naar vi tage Hensyn til Forholdene ved
deres Opstaaen, at forudsætte en sammensat Proces (11).
\end{quote}
\dcite{Høffding, Den menneskelige Tanke, p. 227  [hoeffding/menneskelige-tanke:545]}
% hoeffding/menneskelige-tanke:545  (Den menneskelige Tanke; p. 227)

\begin{quote}
Der viste sig ved Undersøgelsen af Aarsagsbegrebet to Synspunkter, der førte Tanken
i forskellig Retning: \tlg{Ækvivalens} og Kontinuitet. Dersom vi have Ret i, at dette
sidste Begreb er det Karakteristiske for det konkrete Aarsagsforhold, medens det
første Begreb stedse faar en vis abstrakt, metodisk Karakter, saa er det ogsaa
berettiget at hævde den indre Sammenhæng mellem Kausalitet og Udvikling. Vi kunne da
ikke gøre Aarsagsrækken til en absolut Identitetsrække, men hvert Led i den faar, jo
mere vi holde os til Erfaringen, sin bestemte Plads (som i den fremskridende
Forskelsrække).
\end{quote}
\dcite{Høffding, Den menneskelige Tanke, p. 228  [hoeffding/menneskelige-tanke:548]}
% hoeffding/menneskelige-tanke:548  (Den menneskelige Tanke; p. 228)

\begin{quote}
Selvom vi tænkte os, at Alt i den materielle Natur kunde forklares efter denne
Opfattelses Principer, blive Kvaliteterne dog staaende som umiddelbare Fakta, og der
gives ingen Forklaring af, hvorledes de rent kvantitative Forskelle for vore Sanser
kunne fremtræde som kvalitative Forskelle. Det kemiske Produkts Egenskaber kunne
ikke udledes af Elementernes Egenskaber, og naar en Art af fysisk \tlg{Energi} faar sit
\tlg{Ækvivalent} i en anden Art af fysisk \tlg{Energi}, saa har dog \tlg{Ækvivalentet} andre
Egenskaber end den første \tlg{Energiform}. Det sansende Subjekt kan dog ikke skabe disse
Kvalitetsforskelle af Intet. I hvert Tilfælde vilde man derved sætte uløselige
psykologiske Vanskeligheder istedetfor de kemiske og fysiske Vanskeligheder, man
skyder tilside.
\end{quote}
\dcite{Høffding, Den menneskelige Tanke, p. 277+  [hoeffding/menneskelige-tanke:649]}
% hoeffding/menneskelige-tanke:649  (Den menneskelige Tanke; p. 277+)

\begin{quote}
Fra et andet Synspunkt er man kommen til et lignende Resultat. Interessen gaar her
ud paa at udskille alle de Biforestillinger og Hjælpebegreber, med hvilke Emnerne,
det for Tanken Givne, uvilkaarligt eller af metodiske Grunde suppleres. Erkendelsen
skal snarest muligt nærme sig til en ren Beskrivelse af en kontinuerlig Proces,
indenfor hvilken Forskellighederne ere reducerede til et Minimum. Dette opnaas, jo
mere kvantitative Bestemmelser afløse kvalitative, hvorved det bliver muligt at
paavise \tlg{Ækvivalensforhold} mellem de Emner, der afløse hverandre i Erfaringen. Det er
da ikke længer den strænge Videnskabs Opgave at „forklare“; hvem der vil have
„Forklaringer“, henvises til Mytologien eller Metafysiken. Videnskabens Maal er at
give en nøjagtig, metodisk Beskrivelse af alle Forhold og Overgange i Erfaringen.
Hvad nu særligt Aarsagsforholdet angaar, saa udtrykkes Emnernes Afhængighedsforhold
bedre ved Hjælp af det matematiske Funktionsbegreb end ved de ubestemte Udtryk
„Aarsag og Virkning“. Tilmed er der det Uheldige ved enhver Anvendelse af Begreberne
Aarsag og Virkning, at to Emner vilkaarligt tages ud af hele den kontinuerlige
Sammenhæng, i hvilken de forekomme, og stilles i Modsætning til hinanden. Og
Tidsforskellen mellem „Aarsag“ og „Virkning“ har ikke nogen principiel Betydning,
naar Forholdet formuleres som et matematisk Funktionsforhold, udtrykkes i en
Ligning. At Tiden ikke kan „vendes om“, betyder egentligt blot, at de fysiske
Størrelsers Ændringer finde Sted i en bestemt Retning. [note: Richard
Avenarius: Kritik der reinen Erfahrung.1 II., p. 332–335.
— Ernst Mach: Zur Analyse der Empfindungen. p. 168. Die
Mechanik in ihrer Entwickelung.6 p. 259 f.] —
\end{quote}
\dcite{Høffding, Den menneskelige Tanke, pp. 283–284  [hoeffding/menneskelige-tanke:661]}
% hoeffding/menneskelige-tanke:661  (Den menneskelige Tanke; pp. 283–284)

\begin{quote}
120. Man kan derfor ingenlunde sige, at Hume's Problem er blevet løst ved
Opdagelsen af \tlg{Energiens} \tlg{Bestaaen}, der gør det muligt at formulere Forholdet mellem
de forskellige \tlg{Energiformer} som \tlg{Ækvivalens}. \tlg{Ækvivalensforholdet} udelukker hverken
Tidsforskel eller Kvalitetsforskel, ligesaa lidt som Identitetsforholdet behøver at
gøre det (se 69 og 73). Naar A og B ere \tlg{ækvivalente}, er det i en vis Forstand
ligegyldigt, om jeg gaar over fra B til A eller fra A til B. Og dog er det i hvert
enkelt virkeligt Tilfælde enten A, der gaar over til B, eller B, der gaar over til
A, — og denne Forskel kan være en Forskel mellem Liv og Død. I den Carnotske
Sætning (Entropisætningen) træder dette tydeligt frem; her er Retningen netop
afgørende. Varme og Bevægelse ere \tlg{Ækvivalenter}; men det er faktisk vanskeligere at
koncentrere Varmen saaledes, at den igen kan omsættes til ganske det samme Kvantum
Bevægelse end at foretage den omvendte Omsætning. Det er en stor Abstraktion at se
bort fra denne Forskel. Den Carnotske Sætning strider ikke mod Sætningen om
\tlg{Energiens} \tlg{Bestaaen}; den er netop en af dens Forgængere. Men den fører over til et
mere konkret Problem, fører os ind i Naturens virkelige Husholdning, hvor netop
Aarsagsrækkernes Retninger blive af Betydning. Et Exempel fra det menneskelige
Omraade, der svarer til det Carnot'ske fra det fysiske Omraade, frembyder Forholdet
mellem Genialitet, Forbilledlighed, nye friske Indskud indenfor den menneskelige
Udvikling, og Efterligningens og Gentagelsens Herredømme. Mellem Opdagelse og
Efterligning kunde der tænkes et lignende Forhold som mellem Varme og Bevægelse, og
ogsaa her kunde der ligge en Fare i Tendensen til Udligning og
Fordeling. [note: Smlg. G. Tarde: Les lois de l'imitation. p.
412–415.]
\end{quote}
\dcite{Høffding, Den menneskelige Tanke, pp. 285–286  [hoeffding/menneskelige-tanke:665]}
% hoeffding/menneskelige-tanke:665  (Den menneskelige Tanke; pp. 285–286)

\begin{quote}
sit lærde og skarpsindige Værk Identité et Realité 1908 har Meyerson
stillet Entropisætningen i skarp Modsætning til \tlg{Ækvivalensprincipet}: hin vidner om
Tidens Realitet, idet Leddene ikke kunne byttes om, medens denne tenderer til at
ophæve Tiden og tilsidst al Forskel.]. Her anerkendes altsaa, at alle Gaader ikke
ere løste ved Paavisning af \tlg{Ækvivalens}. Der er et historisk Element, som man ikke
kan slippe fra. Det store Spørgsmaal er, i hvilken Grad Aarsagsprocessernes Retning er overladt til dem selv. Ligesom \tlg{Energisætningen} kun kan paavises indenfor
Sammenhæng, der tænkes afsluttede, saa at der ingen Tilførsel eller Afgang af \tlg{Energi}
kan ske, saaledes gaar det ogsaa med Entropisætningen. Ad begge Veje viser det sig,
at det er en Abstraktion, naar vi holde os til de kausale Processer uden at tage
Hensyn til den reale Helhed, indenfor hvilken de foregaa, og som kan indeholde
Tendenser af modsat Retning.
\end{quote}
\dcite{Høffding, Den menneskelige Tanke, p. 287  [hoeffding/menneskelige-tanke:667]}
% hoeffding/menneskelige-tanke:667  (Den menneskelige Tanke; p. 287)

\begin{quote}
Den Analogi, hvorpaa Aarsagssætningen bygges, bliver det vanskeligere at gennemføre,
jo konkretere og individuellere det Emne er, der skal forklares, og det vil sige: jo
mere det udgør en Helhed og gennemløber en Udvikling. Der er her en hel Skala med
stigende Komplikationer fra den simple Stedbevægelse gennem de fysiske og kemiske
Processer til de organiske, psykiske og sociale Fænomener. Jo længere vi komme frem
i denne Række, des mere antager Aarsagsprincipet Karakteren af en ideal Fordring, vi
kun med ringe Grad af Tilnærmelse kunne fyldestgøre. Det bliver stedse vanskeligere
at opløse de sammensatte Processer i simple Aarsagsrækker, — stedse vanskeligere at
opløse de enkelte komplicerede Emner i Elementer, der hvert for sig skulde kunne
betragtes som Udgangspunkter. Stedse vanskeligere bliver det derfor ogsaa at
godtgøre, at det er ganske de samme Emner og de samme Tilstande, vi have for os, naar
en Proces synes at gentage sig. Den gamle Sætning, at Naturen stedse er i
Overensstemmelse med sig selv (sibi consona, som Newton sagde) og ingen Spring gør,
synes ofte at modsiges af Iagttagelsen. Netop den nyeste Tids Forsken yder Exempler
derpaa. I Radium opdagedes et Stof, hvis \tlg{Energiudladninger} kunde synes at stride mod
Sætningen om \tlg{Energiens} \tlg{Bestaaen}. I Modsætning til den ældre Darwinisme hævder man
paa Iagttagelsens Grund den pludselige Opstaaen af nye Typer (ved Mutation), der
foreløbigt ikke kan godtgøres at være Resultatet af en successiv Proces. (Smlgn. dog
92). En Mester i beskrivende Psykologi som William James hævder, at der kan
ske nye Indskud i den sjælelige Udvikling, som ikke kunne udledes af de foregaaende
Processer, og forbereder os idethele paa, at vi ofte maa lade os nøje med en rent
ydre Forbindelse af Emner, en blot „Konjunktion“ uden rationel Forbindelse, et „og“
uden et „derfor“.
\end{quote}
\dcite{Høffding, Den menneskelige Tanke, pp. 287–288  [hoeffding/menneskelige-tanke:669]}
% hoeffding/menneskelige-tanke:669  (Den menneskelige Tanke; pp. 287–288)

\begin{quote}
Naar vi saa ofte stanse i vor Aarsagssøgen paa det psykologiske Omraade, er det især
paa Grund af, at det her er langt vanskeligere end paa andre Omraader at paavise
identiske Forhold og fuldkommen \tlg{Ækvivalens}. Men at dette er Mere end en Gradsforskel,
har Ingen Ret at paastaa.
\end{quote}
\dcite{Høffding, Den menneskelige Tanke, p. 293+  [hoeffding/menneskelige-tanke:681]}
% hoeffding/menneskelige-tanke:681  (Den menneskelige Tanke; p. 293+)

\begin{quote}
127. Erkendelsesproblemet kan kun stilles ud fra det menneskelige Standpunkt. Det er
altsaa et specielt \$S\$, der ligger til Grund. Vi kunne kalde det \$S\_m\$. Det
menneskelige Subjekt har visse kvalitativt udprægede Ejendommeligheder i sin Sansning
og i sin Tænkning, som vi maa tage som rene Fakta (125 α og β). Deres Nødvendighed og
absolute Gyldighed kan ikke godtgøres. Ad Tankeexperimentets Vej kunne vi tænke os et \$S\_n\$ et \$S\_p\$ o. s. v. — subjektive, psykiske \tlg{Energier} med andre specielle
Former, Emner og Motiver end \$S\_m\$ — i Analogi med Matematikernes Hypoteser om andre
Rum end det euklidiske. Desværre kunne vi ikke konstruere disse andre Subjektiviteter; men vi benytte Tankeexperimentet til at indskærpe de rent faktiske Forudsætninger, der
ere givne for en Erkendelse med \$S\_m\$. Man kunde vel tænke sig, at der til de samme optiske Brydningsforskelle svarede andre Synskvaliteter end vore, og man kunde tænke
sig, at de logiske Principer, der staa for os som selvfølgelige, vare afledede af mere
universelle Principer, ligesom Inerti- og \tlg{Energisætningen} maaske ere specielle
Tilfælde af mere almindelige Sætninger.
\end{quote}
\dcite{Høffding, Den menneskelige Tanke, p. 303  [hoeffding/menneskelige-tanke:714]}
% hoeffding/menneskelige-tanke:714  (Den menneskelige Tanke; p. 303)

\begin{quote}
I Afsnittet om Tankens Psykologi blev allerede Forholdet mellem det Aandelige og det
Materielle drøftet, men væsentligst fra et psykologisk, empirisk og metodologisk
Synspunkt. Det gjaldt især om at finde en Synsmaade, der paa en Gang udtrykte, hvad vi virkeligt vide om dette Forhold, og hvad der egnede sig til som Arbejdshypotese at lægges
til Grund ved fortsat Forsken. Hvad vi virkeligt vide, er, at der bestaar et saadant
Forhold mellem de to Arter af Emner, at der kan drages Slutninger fra den ene til den
anden, — at der altsaa bestaar et Funktionsforhold i matematisk Betydning imellem dem.
Derimod var det tvivlsomt, om der ogsaa bestaar et Tidsforhold imellem dem, med
Omsætninger fra fysisk \tlg{Energi} til psykisk, og omvendt. Men hint Funktionsforhold er nu
ogsaa det, der først og fremmest er Brug for under den fortsatte Forsken, saa at det egner
sig til at udtrykke Forholdet mellem de Emner, Psykologien beskæftiger sig med, og dem,
Fysiologien beskæftiger sig med.
\end{quote}
\dcite{Høffding, Den menneskelige Tanke, pp. 329–330  [hoeffding/menneskelige-tanke:769]}
% hoeffding/menneskelige-tanke:769  (Den menneskelige Tanke; pp. 329–330)

\begin{quote}
ved, at der indføres indføres andre Elementer, der kunne drage \tlg{Energien} til sig. En Trang eller en
Drift kan ganske vist dø en naturlig Død af Mangel paa Næring; men en saadan Afdøen vil let
tære paa \tlg{Energien} idethele og volde et afgørende Tab for Sjælehelheden. Hvor Trangen ikke
ophæves ved, at Næringen unddrages, maa der vækkes en anden Trang eller Drift, der kan
hæmme den første. Som rent negativ Dyd er Selvbeherskelse psykologisk umulig [note: Smlg.
mine Etiske Undersøgelser. p. 70.]. Derimod er det muligt, at en Retningsændring af
den i Trangen eller Driften virkende \tlg{Energi} kan ske ved Indflydelse af nye Tilskyndelser,
saa at \tlg{Energien} kan overføres paa andre Elementer. Der maa paa en eller anden Maade skaffes
psykiske \tlg{Ækvivalenter}; det er gennem deres Række, at den personlige Helheds Udvikling
skrider frem. I Platon's Symposion er Ideen om det personlige Livs
kontinuerlige Udvikling gennem psykisk \tlg{Ækvivalens} for første Gang udtalt, og Muligheden af
en Trinrække fra de mest elementære Livsytringer til de højeste (de for Platon højeste)
Aandsformer er paavist. Hvad Platon tænker sig foregaaet i den enkelte Individualitets
Livsløb, foregaar i Virkeligheden stedse indenfor en social Gruppe, hvorfra Tilskyndelser
af forskellig Art og dermed udløsende Kræfter og udfyldende Værdier komme. Det sociale
Synspunkt for det Etiske er ikke blot ved den første Grundlæggelse af et Personlighedsliv,
men ogsaa ved selve Frigørelses- og Selvstændiggørelsesprocessen af afgørende Betydning.
Det er ikke blot det positive, men ogsaa det negative Forhold til Gruppen, der kan føre den
Enkeltes Udvikling fremad. Gennem æggende Paavirkninger kan der udløses Kræfter, der ellers
vilde slumre.
\end{quote}
\dcite{Høffding, Den menneskelige Tanke, pp. 351–352  [hoeffding/menneskelige-tanke:827]}
% hoeffding/menneskelige-tanke:827  (Den menneskelige Tanke; pp. 351–352)

\begin{quote}
160. Troen paa Værdiens \tlg{Bestaaen} kan selv faa Værdi ved at holde Modet oppe under Kampen for
Livet og ved at anspore til at søge og finde nye Værdier som \tlg{Ækvivalenter} for de Værdier, der
forsvinde. Religionen kan paa dette Punkt nærme sig til Etiken i den Grad, at den staar som
en nødvendig Forudsætning for al Handlen. Dette var Fichte's Opfattelse (i hans første
Periode). Han lærte en Identitet af Moral og Religion. Naar jeg sætter mig et etisk Formaal,
maa jeg ogsaa antage, at dets Opnaaelse er mulig, og deri ligger allerede en Tro paa en
moralsk Verdensorden. Der maa ifølge Fichte være en Harmoni mellem Idealernes og Erfaringens
Verden (altsaa mellem Værdi og Virkelighed), dersom etisk Handlen overhovedet skal være
mulig. Fichte's „Ateisme" bestod i, at han gjorde sand Religion til Et med Tro paa en saadan
Harmoni. At denne Tro fandt Udtryk i Antagelsen af et personligt Væsen som Aarsag til
Harmonien, var for ham noget Afledet, der kunde have sine store
Betænkeligheder [note: Fichte: Appellation gegen die Anklage des Atheismus.
1799. p. 34–41.].
\end{quote}
\dcite{Høffding, Den menneskelige Tanke, p. 384  [hoeffding/menneskelige-tanke:881]}
% hoeffding/menneskelige-tanke:881  (Den menneskelige Tanke; p. 384)

\begin{quote}
Paa den anden Side vil der sikkert stedse være Trang til noget Mere end en saadan
Forudsætning. Der vil være en Trang til at finde Udtryk for den samlede Livserfaring i en
Livspoesi, ved hvilken de overleverede Billeder kunne bruges, naar Fantasien ikke paa egen
Haand kan frembringe nye. Poesien udspringer altid af en kraftig Leven, og at Religionerne
have formaaet at give Stødet til en saa stor og mægtig Poesi, hænger nøje sammen med den
koncentrerede Maade, paa hvilken Menneskene i den Religion, der virkeligt er deres Ejendom, føle og leve Livet og bevæges af de Skæbner, det medfører. Ogsaa fra dette Synspunkt set, vil
Religionen ikke kunne forsvinde uden \tlg{Ækvivalent} — hvis der ikke skal ske en Nedgang i det
aandelige Livs Værdi. Hvis man definerer Religion som en Tro paa Værdiens \tlg{Bestaaen}, er det en
religiøs Tro, at de historiske Former for Religion ikke kunne forsvinde uden \tlg{Ækvivalenter}.
\end{quote}
\dcite{Høffding, Den menneskelige Tanke, pp. 384–385  [hoeffding/menneskelige-tanke:883]}
% hoeffding/menneskelige-tanke:883  (Den menneskelige Tanke; pp. 384–385)

\subsection{\textit{Religion og Videnskab} (1910)}

\begin{quote}
Ved Opdagelsen af Radium syntes saaledes Sætningen om \tlg{Energiens} \tlg{Bestaaen} at
rokkes. Her viste sig en stadig Udstraaling af \tlg{Energi}, uden at det kunde
paavises, hvorledes denne \tlg{Energi} var bleven aflejret i dette Stof. Det
stillede da nye Spørgsmaal og gav Anledning til ny Eftertanke. Men det
ændrer Intet i den gamle Metode, hvad det Principielle angaar. — Omtrent
paa samme Tid blev det i Biologien paavist, at Organismer kunne undergaa
pludselige Ændringer, hvorved nye Typer dannes. Derved blev det nødvendigt
at opgive den af Darwin fremsatte Antagelse, at nye Typer opstaa ved, at
smaa Variationer begunstiges og udvikles under Indflydelse af Kampen for
Tilværelsen. Mutationerne, de pludselige Typeændringer, ske nemligt
uafhængigt af Tilværelseskampen; først efterat Mutationen er foregaaet,
bliver Spørgsmaalet, om den nye Type kan bestaa. Men heller ikke her opgives
den Forudsætning, at enhver Begivenhed maa hænge sammen med en anden
Begivenhed. Opgaven bliver at finde de indre og ydre Betingelser, under
hvilke Mutationen foregaar. — Et tredie Exempel kunne vi hente fra
Psykologien, hvor i de sidste Tiaar fra forskellige Sider Opmærksomheden er
bleven henledet paa pludselige Ændringer af Karakter og Personlighed,
Nydannelser paa det psykiske Omraade, især indenfor Følelses- og
Villieslivet. Ogsaa her vise Forholdene sig mere indviklede, end man maaske
en Tid havde trot. Men de principielle Synsmaader vedblive at være de samme;
kun er der større Vanskelighed ved at gennemføre dem.
\end{quote}
\dcite{Høffding, Religion og Videnskab, pp. 25–26  [hoeffding/religion-og-videnskab:33]}
% hoeffding/religion-og-videnskab:33  (Religion og Videnskab; pp. 25–26)

\subsection{\textit{Relation som Kategori} (1921)}

\begin{quote}
Paavisning af kontinuerlige Emnerækker med uombyttelige Led er endnu
ikke Videnskabens højeste Ideal. Den raader vel foreløbig Bod paa det
udvortes Forhold mellem, hvad man kalder Aarsag, og hvad man kalder
Virkning, som fremkaldte Hume's Skepsis; men den gør ikke fuldstændig
Ende derpaa. Et Skridt videre fører Paavisningen af tvesidig \tlg{Ækvivalens}
mellem forskellige Emner. Der kan da være Mulighed for baade at udlede
det ene Emne af det andet og dette andet af det første. Hvor
\tlg{Ækvivalensforholdet} formuleres som en Ligning, fremtræder dette tydelig.
Det er (i Teorien) ligegyldigt, fra hvilken Side af Lighedstegnet man
begynder, og den kvalitative Forskel, saavel som Tidsforskellen mellem
de to Emner har da ikke længere nogen Betydning; et absolut
Identitetsforhold er naaet. Kvaliteterne (til hvilke ogsaa Tiden,
Successionen kan regnes) staar da som blotte Antropomorfismer, subjektive Emner, som den strenge
Videnskab ikke bryder sig om. Man gaar her et Skridt videre end den
mekaniske Naturopfattelse, der ogsaa saa bort fra Kvaliteter i snevrere
Forstand, men endnu anerkendte Tiden som objektiv Form, idet den saa
Bevægelse i Rummet som
\end{quote}
\dcite{Høffding, Relation som Kategori, p. 50+  [hoeffding/relation-som-kategori:95]}
% hoeffding/relation-som-kategori:95  (Relation som Kategori; p. 50+)

\begin{quote}
I Modsætning til en saadan Opfattelse, der ofte hævdes fra matematisk og
naturvidenskabelig Side, og som i filosofisk Form er gennemført af
Marburgerskolen, maa det hævdes, at den Verden af rene Ligninger, i
hvilken streng Videnskab ender, eller som den i hvert Tilfælde stiler hen
imod, ikke mister sin Betydning, fordi det stadig fastholdes, at et
Forhold mellem før og efter aldrig fuldstændig kan gaa op i et Forhold af
ren Identitet. Fordi der er tvesidigt \tlg{Ækvivalensforhold} mellem to
kvalitativt forskellige Emner (som Varme og Bevægelse), er det dog ikke
ligegyldigt, hvilket Emne der kommer først, og hvilket der følger efter.
Der sker noget forskelligt i de to Tilfælde, hvilket tydelig vil vise
sig, naar Rækken føres videre gennem nye Led: selv om der er
\tlg{Ækvivalensforhold} mellem A og B, vil der fra B kunne være en Overgang
mulig til C, som ikke vilde være mulig fra A. Desuden foregaar (hvad
Carnot's Sætning viser med Hensyn til Forholdet mellem Varme og
Bevægelse) Overgangen faktisk ikke lige let i begge Retninger. Men trods
Alt har hin strengt matematiske Afslutning af Naturerkendelsen sin store
Betydning derved, at der til det rent tidløse, logisk-matematiske Forhold
mellem Ligningernes Led svarer den tidslige Over
\end{quote}
\dcite{Høffding, Relation som Kategori, p. 51+  [hoeffding/relation-som-kategori:99]}
% hoeffding/relation-som-kategori:99  (Relation som Kategori; p. 51+)

\begin{quote}
Paa Grund af Sætningen om Bevægelsens Relativitet kunde rent formelt — (matematisk) Himmelfænomenerne lige

saa godt forklares ved, at Solen gik rundt om Jorden, som ved, at Jorden gik
rundt om Solen. Og Tycho Brahe havde opstillet en Mellemform mellem de to
Hypoteser, ifølge hvilken Jorden stod fast, medens Solen og Planeterne gik
rundt om den. Gassendi anstillede en Sammenligning mellem de tre saaledes
opstillede Hypoteser og søgte at vise, at man rent videnskabelig ikke kunde
vælge mellem dem. Men saa paaviste Newton i sit berømte Værk, at de Love,
Kepler havde opdaget for Planetbevægelse, kunde udledes, hvis man udstrakte
Tyngdeloven til at gælde for Himmellegemerne, og da maatte det nødvendigvis
være Jorden og Planeterne, der bevægede sig om Solen. I Kraft af
Faldbevægelsens Lov blev altsaa tilsidst kun én Hypotese mulig. Paa analog
Maade førte omtrent samtidig Leibniz ud over den Skepsis, Bevægelsens
Relativitet tilsyneladende motiverede, idet han hævdede, at det afgørende for
Erkendelsen ikke var selve Bevægelsen, men Bevægelsens Love, og her opstillede
han saa — under Kritik af Descartes' Lære om Bevægelsens \tlg{Bestaaen} —
Sætningen om \tlg{Kraftens} \tlg{Bestaaen}. Her var da en Realitet, der ikke rettede sig
efter Iagttagerens Stade i Rummet. De to store Forskere førte, hver ad sin
Vej, ud over den tilsyneladende Subjektivisme og tog det afgørende Stade i
Lovene for Bevægelsen.
\end{quote}
\dcite{Høffding, Relation som Kategori, pp. 65–66  [hoeffding/relation-som-kategori:129]}
% hoeffding/relation-som-kategori:129  (Relation som Kategori; pp. 65–66)


\end{document}
